\begin{pastexam}[label={px:ntu-111-bime-ds-01}]{臺大生機所 111 學年度（C）第 1 題}
（問答）

(20\%) Given a set of intervals in an array, write a pseudocode function that print all
non-overlapping intervals after merging the overlapping intervals.

Note:

\quad Input example 1: a[0] = (1,3), a[1] = (2,4), your function should print `(1,4)'

\quad Input example 2: a[0] = (1,3), a[1] = (3,5), your function should print `(1,3), (3,5)'
\end{pastexam}

\begin{pastsoln}
先說明解讀：由範例 2，$(1,3)$ 和 $(3,5)$ 只在端點 $3$ 相碰，題目不把它們合併，所以「重疊」指兩個
區間有共同的內部：$(s_1, e_1)$ 和 $(s_2, e_2)$ 在 $s_1 \le s_2$ 時重疊，若且唯若 $s_2 < e_1$。 輸出的
順序題目沒有規定，下面由左到右印。 陣列照資結書的習慣從 $1$ 編號（題目的 \texttt{a[0]} 對應
$A[1]$），每個區間有 $\mathit{start}$、$\mathit{end}$ 兩個欄位，$\mathit{start} < \mathit{end}$。

\textbf{想法。} 先依左端點由小到大排序。 排好之後，會被合併在一起的區間一定是連續的一段：從左到右
掃，維護目前這一群的聯集 $[s, e)$；下一個區間的左端點 $< e$ 就併進來（右端點取大的），否則這一群
已經結束，印出來，從下一個區間開始新的一群。
\begin{pseudo}{合併重疊的區間}
\Procedure{Merge-Intervals}{$A, n$}
    \If{$n = 0$}
        \State \Return
    \EndIf
    \State 依 $\mathit{start}$ 由小到大排序 $A[1..n]$ \Comment{例如合併排序，$O(n \lg n)$}
    \State $s \gets A[1].\mathit{start}$; $e \gets A[1].\mathit{end}$
    \For{$i \gets 2$ \To $n$}
        \If{$A[i].\mathit{start} < e$} \Comment{和目前這一群重疊}
            \State $e \gets \max(e, A[i].\mathit{end})$
        \Else
            \State 印出 $(s, e)$ \Comment{這一群結束}
            \State $s \gets A[i].\mathit{start}$; $e \gets A[i].\mathit{end}$
        \EndIf
    \EndFor
    \State 印出 $(s, e)$
\EndProcedure
\end{pseudo}

\textbf{為什麼對。} 迴圈不變量：處理完前 $i - 1$ 個（排序後的）區間時，(1) 已經印出的區間就是前面
幾群各自的聯集，彼此不重疊；(2) $(s, e)$ 是目前這一群（包含 $A[i-1]$）的聯集，$s$ 是這群最小的左端點，
$e$ 是最大的右端點，而且這群裡的區間經由重疊串成一片，聯集是一整段。 看 $A[i]$：
\begin{itemize}
\item $A[i].\mathit{start} < e$：$A[i]$ 和這一群聯集的內部有交集（$A[i].\mathit{start} \ge s$，因為排過
序），應該併進來；新的聯集是 $(s, \max(e, A[i].\mathit{end}))$。
\item $A[i].\mathit{start} \ge e$：$A[i]$ 和這一群不重疊。 之後的區間左端點都 $\ge A[i].\mathit{start}
\ge e$，也都不會和這一群重疊，所以這一群已經完整，可以印出。
\end{itemize}
迴圈結束後最後一群還沒印，所以最後再印一次。 每一群印出時都已完整，而且相鄰兩群之間有空隙或只碰
到端點，輸出的區間彼此不重疊。

\textbf{範例。} 範例 1：排序後 $(1,3), (2,4)$；$2 < 3$，併成 $(1,4)$，印出 $(1,4)$。 範例 2：$(1,3),
(3,5)$；$3 < 3$ 不成立，印出 $(1,3)$，再印出 $(3,5)$。 再一例：$(6,8), (1,9), (2,4), (4,7)$ 排序成
$(1,9), (2,4), (4,7), (6,8)$，後面三個的左端點都 $< 9$，全部併進來，印出 $(1,9)$。 這個例子說明
$e$ 要取 $\max$：只看相鄰兩個區間會把 $(2,4)$、$(4,7)$ 誤判成不相連。

\textbf{複雜度。} 排序 $O(n \log n)$，掃描 $O(n)$，總共 $O(n \log n)$；除了排序外只用 $O(1)$ 的額外空間。
如果題目要把端點相碰也算重疊，把條件改成 $A[i].\mathit{start} \le e$ 即可。
\end{pastsoln}
